% HD-N19. GPT-6 Astra Pro /root/imprimitive. 21 November 2026.
% Inherited effort; the numerical setting is not exposed.
\section{HD-N19: the rank-two Atiyah quotient and the index-three Euler obstruction}
\label{sec:rank-two-atiyah}

\paragraph{The actual objects and the conclusion.}
Let $X$ be an abelian sixfold and let $E$ be an actual locally free
bundle of rank two. Suppose that
\[
 c_1(E)=0,\qquad c_2(E)=d\eta^2,\qquad d>0,                 \tag{N19.1}
\]
where $\eta\in H^{1,1}(X,\mathbb R)$ is nondegenerate and its
Hermitian form has three positive and three negative eigenvalues.
No existence assertion for such a bundle is made here. Every bundle
in this explicitly specified class has a nonzero kernel in at least
one of its degree-zero and degree-two categorical traces. In particular,
every simple bundle in this class fails ordinary semiregularity.
This is semiregularity in Perry's Definition2.4, namely injectivity
in degree two; no degree-four trace hypothesis is added.

The proof gives the missing degree-four map explicitly. For every
bundle(N19.1), the actual Hochschild evaluation in degree four has
image of dimension exactly30. Thus injective degree-two trace forces
$\operatorname{ext}^2(E,E)=30$ by the sixfold Serre pairing.
The last paragraph records the exact constraints for arbitrary
coherent finite point actions and the stronger ordinary-descent result.

\paragraph{1. The global curvature identity and its contraction.}
The integral group $H^2(X,\mathbb Z)$ is torsion free. Thus the first
equality in(N19.1), even when stated rationally, puts $\det E$ in
$\operatorname{Pic}^0(X)$. Give this line bundle its unitary flat metric.
Starting with any Hermitian metric on $E$, multiply it by the positive
square root of the ratio of determinant metrics. The resulting Chern
connection induces the chosen flat determinant connection, so its
curvature is trace free.

Use Chern-normalized Dolbeault and HKR conventions, in which normalized
curvature $\mathcal F$ represents the Atiyah class used in
$\operatorname{ch}(E)=\operatorname{tr}\exp(\mathcal F)$ and its contractions
are the degree-one Hochschild evaluations. With raw curvature this
means the explicit scalar change of tangent generators by $1/(2\pi i)$,
with the common sign changed when the Atiyah class is defined with the
opposite sign. These invertible scalar changes do not alter any
evaluation image. The precise curvature comparison is Kapranov,
\emph{Rozansky--Witten invariants via Atiyah classes},
Compositio Math.115 (1999),71--113, Propositions1.3.1 and1.3.2
on pp.75--76, followed by(1.4.1) onp.76;
\href{https://doi.org/10.1023/A:1000664527238}
{DOI10.1023/A:1000664527238}.

Locally $\mathcal F$ is a trace-free two-by-two matrix whose entries
are differential forms of even total degree two. Its entries commute
under wedge product. The literal matrix Cayley--Hamilton identity gives
\[
 \mathcal F^2=\frac12\operatorname{tr}(\mathcal F^2)
                        \operatorname{id}_E.                  \tag{N19.2}
\]
This is invariant under every change of bundle frame, hence is a
global identity of endomorphism-valued forms. Its scalar class is
\[
 \left[\frac12\operatorname{tr}(\mathcal F^2)\right]
       =\operatorname{ch}_2(E)=-c_2(E)=-d\eta^2
                      \quad\text{in }H^2(X,\Omega_X^2).         \tag{N19.3}
\]
The Dolbeault equality follows from Hodge comparison on compact
K\"ahler $X$. The scalar difference is $\bar\partial$-exact;
multiplying its primitive by $\operatorname{id}_E$ gives the
corresponding equality in the endomorphism Dolbeault complex.

Choose invariant holomorphic one-forms $z_1,\ldots,z_6$, with dual
vectors $v_i$, and write the invariant representative
$\eta=\sum_i z_i\wedge b_i$. Nondegeneracy says precisely that
$b_1,\ldots,b_6$ is a basis of $H^1(X,\mathcal O_X)$.
Write $\mathcal F=\sum_i z_i\wedge f_i$ and set
\[
 A_i=[f_i]\in\operatorname{Ext}^1(E,E),\qquad
 B_i=[b_i]\operatorname{id}_E\in\operatorname{Ext}^1(E,E).
\]
These are the actual evaluations of the twelve chosen Hochschild
generators. Their products anticommute: evaluation is the characteristic
ring homomorphism into the graded centre. The precise source is
Buchweitz--Flenner, \emph{The global decomposition theorem for Hochschild
(co-)homology of singular spaces via the Atiyah--Chern character},
Adv. Math.217 (2008),243--281, Section6.3, Proposition6.3.1 and
Corollary6.3.2; author version
\href{https://arxiv.org/pdf/math/0606730}{arXiv:math/0606730},
pp.30--31. On $X$ the Hochschild ring is the exterior algebra on
$H^1(\mathcal O_X)\oplus H^0(T_X)$. Its multiplicative identification
can be read directly from the group-difference isomorphism
$(x,y)\mapsto(x,y-x)$: the diagonal becomes $X\times\{0\}$,
whose formal Koszul resolution gives the exterior Yoneda algebra,
tensored with $H^*(\mathcal O_X)$.

The contraction sign is retained. For $i<j$, the coefficient of
$z_i\wedge z_j$ in $\mathcal F^2$ is $-(f_if_j-f_jf_i)$, with
cohomology class $-2A_iA_j$. The corresponding coefficient of
$\eta^2$ is $-2b_ib_j$. Equations(N19.2)--(N19.3) give
\[
                         A_iA_j=-dB_iB_j\qquad(i\ne j).         \tag{N19.4}
\]
Contraction with invariant vectors commutes with $\bar\partial$.
Thus this passage to Yoneda classes does not require a global
trivialization of $E$.

\paragraph{2. The actual exterior quotient in degrees at least three.}
Put $s^2=-d$ over $\mathbb C$ and consider the explicit quotient
\[
 R=\frac{\bigwedge\langle A_1,\ldots,A_6,B_1,\ldots,B_6\rangle}
             {(A_iA_j-s^2B_iB_j:i<j)}.                         \tag{N19.5}
\]
Evaluation factors through $R$. Every monomial reduces, by replacing
a pair of $A$'s, to a monomial with no $A$ or one $A$.
For $k\ge3$, a remaining monomial $A_iB_J$ with $i\in J$ is zero:
choose $j\in J\setminus\{i\}$ and multiply
$A_iA_j=s^2B_iB_j$ by $A_i$ to obtain $s^2A_iB_iB_j=0$.
If $i\notin J$, all indices are distinct. For three distinct
indices, the three reductions of $A_iA_jA_k$ give
\[
                A_iB_jB_k=A_kB_iB_j=-A_jB_iB_k.                \tag{N19.6}
\]
Multiplying by further $B$'s identifies all choices of the single
$A$ index in a fixed $k$-element index set, with the exterior signs.
There are therefore at most two spanning elements for each such set:
\[
 \dim R_k\le2\binom6k\ (3\le k\le6),\qquad R_k=0\ (k>6).        \tag{N19.7}
\]
There are no further relations in this spanning set. The two
homomorphisms $R\to\bigwedge\langle b_1,\ldots,b_6\rangle$ given by
$(A_i,B_i)\mapsto(\pm s b_i,b_i)$ carry the two elements of a
fixed index set to $(b_I,b_I)$ and $(s b_I,-s b_I)$ in their
direct sum. Those columns are independent. Thus equality holds
in(N19.7), and in particular
\[
             \dim\operatorname{im}\operatorname{ev}_E^4
                           \le\dim R_4=30.                     \tag{N19.8}
\]
This bounds the actual evaluation image. It does not assert that
every self-Ext class is generated in degree one.

\paragraph{3. Character ranks and the Serre adjoint.}
The rank-two Chern-root identity, with $c_1=0$ and $c_j=0$ for
$j>2$, gives the whole character
\[
 \operatorname{ch}(E)=e^{s\eta}+e^{-s\eta}
   =2-d\eta^2+\frac{d^2}{12}\eta^4-\frac{d^3}{360}\eta^6.        \tag{N19.9}
\]
The trace/evaluation identity identifies
$c_E^k:=\sigma_E^k\operatorname{ev}_E^k$ with the actual
degree-$k$ Hochschild action on(N19.9). The exact published
comparison is Buchweitz--Flenner, Proposition6.2.1 and Theorem6.4.2,
loc. cit., author-version pp.29--31.

Set $U=H^{0,1}(X)\oplus H^{1,0}(X)^*$ and
$\pi_\pm(\alpha,v)=\alpha\pm s\iota_v\eta$.
The map $(\pi_+,\pi_-):U\to H^{0,1}\oplus H^{0,1}$ is an
isomorphism because $s\ne0$ and $\eta$ is nondegenerate.
Its two graph subspaces separate the actions on the two
exponentials. A product using both graph subspaces kills both
exponentials. The two degree-four images are
\[
       \bigwedge^4H^{0,1}\,e^{s\eta},\qquad
       \bigwedge^4H^{0,1}\,e^{-s\eta}.                          \tag{N19.10}
\]
They intersect trivially. An equality first identifies the lowest
components as $\xi\in H^{0,4}$; the next components give
$2s\eta\wedge\xi=0$. The map
$\eta\wedge:H^{0,4}\to H^{1,5}$ is injective: in
$\eta=\sum z_i b_i$, its coefficient at $z_i$ is $b_i\wedge\xi$,
and all six vanish only when every coefficient of $\xi$ vanishes.
The same argument in degree one gives
\[
                 \operatorname{rank}c_E^4=30,\qquad
                 \operatorname{rank}c_E^1=12.                 \tag{N19.11}
\]
In particular(N19.8) is an equality.

Fix a nonzero invariant holomorphic volume form. Calabi--Yau Serre
duality identifies the dual of $\sigma_E^2$ with
$\operatorname{ev}_E^4$, after the volume-form identification of the
relevant Hochschild homology dual with $HH^4(X)$. This is the
literal pairing used in HD-N11: cyclicity identifies it with
$\operatorname{tr}(\xi\,\operatorname{ev}_E^4(a))$ for
$\xi\in\operatorname{Ext}^2(E,E)$ and $a\in HH^4(X)$.
An injective $\sigma_E^2$ therefore makes
$\operatorname{ev}_E^4$ surjective. Equations(N19.8) and(N19.11) give
\[
 \operatorname{ext}^4(E,E)=\operatorname{ext}^2(E,E)=30,\qquad
                        \operatorname{ext}^1(E,E)\ge12.        \tag{N19.12}
\]
The quotient upper bound is essential here. No inference from
injectivity of $\sigma^2$ to injectivity of $\sigma^4$ was used.

\paragraph{4. Degree-zero trace and the rank-two splitting map.}
Degree-zero categorical trace kills commutators and nilpotents.
For bundle endomorphisms, the latter statement follows from the
actual kernel filtration, whose successive induced maps are zero,
and additivity of every Atiyah-trace component. Generalized-trace
additivity with its coherence hypothesis is
Carmeli--Cnossen--Ramzi--Yanovski, Forum Math. Sigma13 (2025),e93,
Theorem5.27 onp.58 and Remark5.25 onp.57,
\href{https://doi.org/10.1017/fms.2025.23}{DOI10.1017/fms.2025.23}.
The actual kernel/cofiber diagrams were checked in HD-R04.
Injective $\sigma_E^0$ thus makes $\operatorname{End}(E)$ a finite
commutative reduced $\mathbb C$-algebra, hence $\mathbb C^r$.
Its idempotents split $E$ into actual nonzero locally free summands,
so $r\le2$.

In fact $r=2$ is impossible. Such a splitting would be $L\oplus M$
with $c_1(M)=-c_1(L)$. For $D=c_1(L)$ this gives
\[
                              D^2=-d\eta^2.                    \tag{N19.13}
\]
Write $D$ and $\eta$ as matrices of $(1,1)$-forms in holomorphic
and antiholomorphic bases. Let $M$ be the former matrix and $N$
the latter, which is invertible. The coefficients in(N19.13)
are the two-by-two minors, so
\[
                    \bigwedge^2(N^{-1}M)=-d\operatorname{id}.   \tag{N19.14}
\]
The operator $T=N^{-1}M$ is invertible and preserves every two-plane:
every decomposable nonzero bivector is carried to the same line.
Intersecting two different two-planes through a fixed line shows
that $T$ preserves every line. Applying this to a basis and its
pairwise sums gives $T=\lambda\operatorname{id}$, hence
$D=\lambda\eta$. Both classes are real, so $\lambda$ is real;
(N19.14) gives $\lambda^2=-d$, a contradiction. Consequently
\[
 \sigma_E^0\text{ injective}\ \Longrightarrow\
                       \operatorname{End}(E)=\mathbb C.        \tag{N19.15}
\]
Conversely simplicity makes $\sigma^0$ injective because the
identity has the nonzero rank-two character.

\paragraph{5. The integral Euler contradiction.}
Put $V_\eta=\int_X\eta^6/6!$. The index-three assumption gives
$V_\eta<0$ by the product of the six Hermitian eigenvalues in
a positive volume basis. Riemann--Roch has Todd class1, so
\[
 \chi(X,E)=-2d^3V_\eta\in\mathbb Z_{>0},\qquad
                    \chi(E,E)=64\chi(X,E)\ge64.                \tag{N19.16}
\]
The factor64 is exact: $E^\vee$ has the same character and the
degree-six coefficient of
$(e^{s\eta}+e^{-s\eta})^2=e^{2s\eta}+2+e^{-2s\eta}$
is $2^6$ times that in(N19.9).

If both traces were injective, (N19.12),(N19.15) and Serre duality
would instead give
\[
 \begin{split}
 \chi(E,E)&=2\operatorname{ext}^0(E,E)-2\operatorname{ext}^1(E,E)
                +2\operatorname{ext}^2(E,E)-\operatorname{ext}^3(E,E)\\
          &\le2-24+60=38,
 \end{split}                                                   \tag{N19.17}
\]
contradicting(N19.16). The same proof starts with
$\operatorname{End}(E)=\mathbb C$ for a simple bundle. This proves
the ordinary obstruction.

\paragraph{6. Exact finite-action and descent scope.}
Let $G\subset X\times\widehat X$ be a faithful finite point subgroup
with an actual coherent translation/tensor action and an equivariant
structure on $E$. The HD-N12 comparison works in dimension six by
the same proof: all nonidentity Hochschild sectors vanish, the
induction/forgetful composite is multiplication by $|G|$, and the
invariant category has the invariant sixfold Serre pairing.
Thus weak $G$-semiregularity in Perry Definition2.6 is precisely
injectivity of the ambient trace on $\operatorname{Ext}^2(E,E)^G$.
Perry's Definition2.6 and Lemma2.7 are on pp.5--6 of the supplied
author version arXiv:2604.00511v2. They impose no degree-zero
trace condition.

All evaluations are invariant. Hence(N19.4)--(N19.12) give, for
every weakly $G$-semiregular bundle in this class,
\[
 e_2^G=e_4^G=30,\qquad e_1^G\ge12,\qquad
               \chi_G(E,E)=\frac{64\chi(X,E)}{|G|}.             \tag{N19.18}
\]
The last equality is the actual eigen-Euler descent computation
of HD-N12. It does not assert descent as an ordinary bundle.
When the degree-zero equivariant trace is injective as well,
the idempotents and(N19.13)--(N19.14) give $e_0^G=1$.
The invariant middle Serre pairing is skew symmetric in degree3,
so $e_3^G$ is even. Every such action satisfies
\[
 \frac{64\chi(X,E)}{|G|}\in\{2,4,\ldots,38\},\qquad
                       |G|\ge\frac{32}{19}\chi(X,E).           \tag{N19.19}
\]
These exact constraints do not exclude every possible finite action.

Two actual stabilizer maps sharpen this scope. For $(a,P)\in G$,
taking determinants of $T_a^*E\otimes P\simeq E$ gives
$P^2\simeq\mathcal O_X$, since the flat determinant line is
translation invariant. There is no nonidentity element with $a=0$.
Its nontrivial $P$ would have order two. The given coherent cyclic
action, after a constant rescaling, makes $E$ a module over
$\mathcal O_X\oplus P$, the algebra of the connected etale double
cover $p:X'\to X$. Explicitly, for $X=V/\Lambda$, the torsion
line $P$ defines a nontrivial character $\Lambda\to\{1,-1\}$;
put $\Lambda'$ equal to its kernel. Then
$p:V/\Lambda'\to V/\Lambda$, and its deck involution is
translation by the class of any element of $\Lambda\setminus\Lambda'$.
Thus $E=p_*L$ for an actual line bundle $L$:
local freeness follows from the finite etale algebra, and connectedness
makes the module rank constant and equal to one. The actual pullback
is $p^*E=L\oplus\tau^*L$. Deck translations act trivially on cohomology,
so $c_1(p^*E)=0$ gives $c_1(L)=0$ and then $p^*c_2(E)=0$.
But $p_*p^*=2$ makes rational pullback injective, contradicting
$d\eta^2\ne0$. Therefore $G$ projects injectively to translations,
and its tensor coordinates lie in $\widehat X[2]$.

For a finite translation group with genuine ordinary descent data,
let $p:X\to Y=X/G$ and $E=p^*E_Y$. Rational pullback is an
isomorphism onto translation invariants, which are all cohomology,
and it carries Hodge structures and the form by its invertible
derivative. Thus $E_Y$ has(N19.1) with the same index three.
The invariant Ext and trace diagrams are the literal descent and
pullback diagrams; also $\chi(X,E)=|G|\chi(Y,E_Y)$. The ordinary
contradiction therefore excludes simultaneous degree-zero trace
injection and weak semiregularity for this genuine descent.
Scalar-twisted category actions and graphs with nontrivial tensor
lines have not been replaced by ordinary descent. For those actions
the proved statements are exactly(N19.18)--(N19.19) and the two
stabilizer constraints above.

\paragraph{Proof and arithmetic record.}
The universal proof uses global curvature, explicit Yoneda relations,
the actual trace/evaluation adjunction, integral Euler characteristics
and the stated descent maps. The companion
\texttt{rank\_two\_atiyah\_quotient\_certificate.py} builds the binomial
exterior ideal over signed powers of $d$. Every consistency cycle has
exponent zero, so its Hilbert vector
\[
                     (1,12,51,40,30,12,2,0,\ldots,0)
\]
has no exceptional specialization at nonzero $d$.
A separate exterior implementation checks48 finite-field character
ranks for an index-three diagonal form and the exact Euler ratio64.
Degree-six character rank is1 whereas degree-six quotient dimension
is2; those two maps are not identified. The proof uses degree four,
where the quotient bound and character rank both equal30.

